\documentclass[10pt,a4paper]{amsart}
\usepackage[margin=19mm]{geometry}
\usepackage{amsmath,amssymb,mathtools}
\usepackage[scr=rsfs,cal=euler]{mathalfa}
\usepackage{xfrac}
\usepackage[unicode,hidelinks,bookmarks=false]{hyperref}
\newcommand{\ie}{i.e.\ }
\newcommand{\cf}{cf.\ }
\newcommand{\resp}{respectively\ }
\newcommand{\Subsection}{\subsection}
\providecommand{\nobreakdash}{\nobreak\hbox{-}}
\setlength{\emergencystretch}{2em}
\multlinegap0pt
\usepackage[shortlabels]{enumitem}
\usepackage{amssymb}
\usepackage{xfrac}
\usepackage{mathtools}
\newtheorem{definition}{Definition}
\newtheorem{theorem}{Theorem}
\DeclareMathOperator{\Cay}{Cay}
\DeclareMathOperator{\Ind}{Ind}
\newcommand{\Z}{\mathbb{Z}}
\newcommand{\quot}[2]{\sfrac{#1}{#2}}
\title{The Algebraicity Problem for Hard-Core Entropy Constants on the Discrete Hypertori}
\author{Yotam Svoray}
\date{Source: arXiv:2608.28332v2 (CC BY 4.0)}
\begin{document}
\maketitle

\noindent\emph{Source and licence.} The Algebraicity Problem for Hard-Core Entropy Constants on the Discrete Hypertori. Version arXiv:2608.28332v2 (CC BY 4.0), distributed by arXiv under the Creative Commons Attribution 4.0 International licence, http://creativecommons.org/licenses/by/4.0/, as recorded in the arXiv licence metadata for this exact version and shown on the abstract page https://arxiv.org/abs/2608.28332v2. Full licence notice: \texttt{licenses/arxiv-2608.28332-CC-BY-4.0.txt}.\par\smallskip

\noindent\textit{Editorial scope.} Counted unit: the article's theorem thm:existence (source0.tex lines 375--383). The article's complete original proof of it is delivered with it (source0.tex lines 384--424, one \texttt{proof} environment of 41 lines). No mathematics is written, repaired or re-derived: the statement and the proof are the article's own lines, in the article's own notation, and no retained line is substituted or re-flowed.  Retained uncounted as the source-local context the proof needs: lines 345--351.  Nothing was omitted from any retained span, so the package carries no omission record.  No retained line contains a citation, so the wrapper carries no bibliography and no bibliography entry.  No retained line contains a figure environment, an image-inclusion command or drawing code, so no image is delivered and the package declares no asset dependency.  Editorial formatting only, in addition: an AMS \texttt{amsart} preamble that restates the article's own declarations and macros used by the retained lines, namely the environments \texttt{definition}, \texttt{theorem} on separate counters, together with the standard packages those lines need.  The retained environments print in this document as Definition 1, Theorem 1 in order of appearance, whereas the article prints the same items with the numbering of its own full document.  The complete native source of the article as distributed in the arXiv e-print for this exact version is bundled unchanged as \texttt{sources/208783/source0.tex}.
\begin{definition}\label{def:torus}
For $d,n\ge1$, the \textbf{$d$-dimensional discrete torus of side $n$} is
\[
T_d(n)\ =\ \Cay\big((\quot{\Z}{n\Z})^d,\,S_n\big),\qquad S_n=\{\pm e_1,\dots,\pm e_d\}\setminus\{0\}\subseteq(\quot{\Z}{n\Z})^d,
\]
where $e_i$ is the $i$-th standard basis vector. We write $a_d(n)=i(T_d(n))$ and we define $\kappa_d=\lim_{n\to\infty}a_d(n)^{1/n^d}$. By convention, we set both $T_d(1)$ and $T_0(n)$ to the single-vertex graphs with no edges, and we set $a_0(n)=2$ for every $n$.
\end{definition}

\begin{theorem}\label{thm:existence}
For every $d\ge1$ the constant $\kappa_d$ is well defined and for every $n\ge1$ we have that 

\begin{enumerate}
    \item $\kappa_d^{\,n^d}\,2^{-dn^{d-1}}\ \le\ a_d(n)\ \le\ \kappa_d^{\,(n+1)^d}$, 
    \item $
\left|\frac1{n^d}\log a_d(n)-\log\kappa_d\right|\ \le\ \frac{d\log2}n.$
\end{enumerate}
\end{theorem}

\begin{proof}
As the proof of this theorem is long, we break it up into steps:\\

\underline{Step 1: edge deletion.} For any graph $G$ and a single edge $uv$, deleting it to form $H=G-uv$ gives that $i(G)\le i(H)\le2i(G)$. This is true since every $G$-independent set is $H$-independent, and every $H$-independent set either avoids one of $u,v$ (and therefore is already $G$-independent) or contains both, in which case removing $v$ leaves a $G$-independent set, and this is injective on such sets, so $i(H)\le i(G)+i(G)$. Iterating over $m$ deleted edges gives $i(G)\le i(H)\le2^mi(G)$. Applied to $G =T_d(n)$ and $H=B_d(n)$, in which $dn^{d-1}$ edges were deleted, we get that:
\begin{equation*}\label{eq:torusgridcompare}
a_d(n)\ \le\ b_d(n)\ \le\ 2^{dn^{d-1}}a_d(n).
\end{equation*}
In particular, we get that $a_d(n)^{\frac{1}{n^d}} \leq b_d(n)^{\frac{1}{n^d}} \leq 2^{\frac{d}{n}}a_d(n)^{\frac{1}{n^d}}$.\\

\underline{Step 2: monotonicity of $b_d$.} For every graph $G$ and every vertex $v$ of $G$, partitioning the independent sets of $G$ by whether they contain $v$ gives the formula $i(G)=i(G-v)+i(G-N[v])\ge i(G-v)$, where $N[v]$ is the closed neighborhood of $v$. Iterating over several deleted vertices shows that deleting vertices from a graph (i.e.\ passing to an induced subgraph) never increases the number of independent sets. Therefore, for $m\le n$, the sub-box $\{0,\dots,m-1\}^d\subseteq\{0,\dots,n-1\}^d$ induces, inside $B_d(n)$, exactly the graph $B_d(m)$ (as the grid-edge relation is local and does not see the ambient box size), so $B_d(m)$ is an induced subgraph of $B_d(n)$ and we get that $b_d(m)\ \le\ b_d(n)$ for every $2 \leq m \leq n$. \\

\underline{Step 3: submultiplicativity along multiples.} Partitioning $\{0,\dots,kn-1\}^d$ into $k^d$ translated copies of $\{0,\dots,n-1\}^d$ and restricting an independent set of the big box to each copy gives an injection $\Ind(B_d(kn))\hookrightarrow\Ind(B_d(n))^{k^d}$ (a $B_d(kn)$-independent set restricted to any sub-box is $B_d(n)$-independent, and the restrictions to the $k^d$ sub-boxes together determine the original set), so $b_d(kn)\ \le\ b_d(n)^{k^d}$.\\

\underline{Step 4: super-multiplicativity via buffering.} Inside $\{0,\dots,k(n+1)-2\}^d$, insert a single empty coordinate layer between each pair of axis-adjacent $n$-blocks (i.e., $k-1$ buffer layers along each axis, so the total side length is $kn+(k-1)=k(n+1)-1$), leaving every buffer vertex unoccupied and placing an arbitrary $B_d(n)$-independent set independently in each of the $k^d$ blocks produces a $B_d\big(k(n+1)-1\big)$-independent set, injectively in the $k^d$-tuple of choices (distinct blocks are separated by an empty layer along every axis, so no adjacency is ever created between blocks). Thus we get that $b_d(n)^{k^d}\ \le\ b_d\big(k(n+1)-1\big)$.\\

\underline{Step 5: existence of the limit.} In order to show that $\kappa_d$ is well defined, we set it to be $\kappa_d=\inf_{n\ge1}b_d(n)^{1/n^d}\in(0,2]$ (which is finite and positive since $1\le b_d(n)\le 2^{n^d}$ term-wise with respect to $n$) and we show that it equals the limit in  Definition~\ref{def:torus}. By the definition of the infimum we get that $b_d(n)\ \ge\ \kappa_d^{\,n^d}$ for every $n \geq 1$, and so $\liminf_n b_d(n)^{1/n^d}\ge\kappa_d$. For the reverse inequality, fix $m\ge2$ and for every $n\ge m$, by the Euclidean algorithm we can write $n=q_n m+r_n$ where $q_n\ge1$ and $0\le r_n<m$, and so $n\le(q_n+1)m$. By the monotonicity of $\{b_d(n)\}_{n=1}^\infty$ and from the fact that $b_d(kn)\ \le\ b_d(n)^{k^d}$, we get that 
\[
b_d(n)\ \le\ b_d\big((q_n+1)m\big)\ \le\ b_d(m)^{(q_n+1)^d}.
\]
By taking $n^d$-th roots we get that $b_d(n)^{1/n^d}\le \big[b_d(m)^{1/m^d}\big]^{(\frac{(q_n+1)m}{n})^d}$. As $n\to\infty$ with $m$ fixed we have that $q_n\to\infty$ and $1\le(q_n+1)\frac{m}{n}<(q_n+1)\frac{m}{qm}=1+\frac{1}{q_n}\to1$, so $(q_n+1)\frac{m}{n}\to1$. By continuity of the function $x\mapsto x^\alpha$ at $\alpha=1$ (with base $b_d(m)^{\frac{1}{m^d}}>0$ fixed) we get that $\limsup_n b_d(n)^{\frac{1}{n^d}}\le b_d(m)^{\frac{1}{m^d}}$. Since $2\le m$ was arbitrary, we have that $\limsup_n b_d(n)^{\frac{1}{n^d}}\le\inf_{m\ge2}b_d(m)^{\frac{1}{m^d}}=\kappa_d$. Note that the infimum over $m\ge2$ agrees with that over all $m\ge1$, since we have that $b_d(m) \leq 2^{m^d}$ and so $b_d(m)^{\frac{1}{m^d}} \leq 2 =b_d(1)$. Combined with the liminf bound, we get that $b_d(n)^{\frac{1}{n^d}}\to\kappa_d$ as $n\to\infty$.\\ 

\underline{Step 6: the limit of $b_d(n)^{\frac{1}{n^d}}$.} Since $b_d(n)^{k^d}\ \le\ b_d\big(k(n+1)-1\big)$, then by sending $k \to \infty$ and writing $N=k(n+1)-1$, we get that
\[
b_d(n)\ \le\ b_d(N)^{\frac{1}{k^d}}\ =\ \Big[b_d(N)^{\frac{1}{N^d}}\Big]^{\frac{N^d}{k^d}}.
\]
Together with Step 5, we can conclude that $b_d(N)^{\frac{1}{N^d}}\to\kappa_d$ as $k\to\infty$ (equivalently $N\to\infty$), and $\frac{N^d}{k^d}=\frac{\big(k(n+1)-1\big)^d}{k^d}\to(n+1)^d$, since $b_d(n)$ is a fixed quantity bounded above by a sequence converging to $\kappa_d^{(n+1)^d}$, we conclude $b_d(n)\le\kappa_d^{(n+1)^d}$. Together with the fact that $a_d(n)\ \le\ b_d(n)\ \le\ 2^{dn^{d-1}}a_d(n)$ and with the fact that $b_d(n)\ \ge\ \kappa_d^{\,n^d}$, we get that $a_d(n)\le b_d(n)\le\kappa_d^{(n+1)^d}$ and $a_d(n)\ge b_d(n)2^{-dn^{d-1}}\ge\kappa_d^{n^d}2^{-dn^{d-1}}$, which is exactly item 1. Taking logarithms and dividing by $n^d$ also shows that $a_d(n)^{1/n^d}\to\kappa_d$. Since $a_d(n)\ \le\ b_d(n)\ \le\ 2^{dn^{d-1}}a_d(n)$, we get that $|\log a_d(n)-\log b_d(n)|\le dn^{d-1}\log2$, and so $\big|\log a_d(n)/n^d-\log b_d(n)/n^d\big|\le d\log2/n\to0$, and $\log b_d(n)/n^d\to\log\kappa_d$ by Step 5.\\

\underline{Step 7: the rate (item 2).} Fix $n\ge1$, for $k\ge1$ we can tile $\{0,\dots,kn-1\}^d$ by $k^d$ copies of $\{0,\dots,n-1\}^d$ and delete every cross-tile edge of $T_d(kn)$, including its $d$ wraparound bundles. A direct count (each of the $d$ axes contributes $k$ inter-tile boundaries, each carrying $(kn)^{d-1}$ edges) shows that $dk(kn)^{d-1}=dk^dn^{d-1}$ edges are deleted, and what remains is the disjoint union of $k^d$ copies of $B_d(n)$ (within a tile, and ignoring the edges leaving it, $T_d(kn)$ restricts to exactly the box grid $B_d(n)$). By the edge-deletion bound of Step 1 (applied with $G=T_d(kn)$, the graph with more edges, and $H$ the disjoint union), we get that
\[
0\ \le\ \log b_d(n)^{k^d}-\log a_d(kn)\ \le\ m\log2,
\]
and by dividing by $(kn)^d=k^dn^d$ and letting $k\to\infty$ (using the fact that $a_d(kn)^{1/(kn)^d}\to\kappa_d$ from Step 6, applied along the subsequence $kn$), we get the one-sided bound
\begin{equation*}
0\ \le\ \frac{\log b_d(n)}{n^d}-\log\kappa_d\ \le\ \frac{d\log2}n.
\end{equation*}
Separately, the inequality $a_d(n)\ \le\ b_d(n)\ \le\ 2^{dn^{d-1}}a_d(n)$ (which comes from deleting $dn^{d-1}$ edges) gives us that
\begin{equation*}
-\frac{d\log2}n\ \le\ \frac{\log a_d(n)}{n^d}-\frac{\log b_d(n)}{n^d}\ \le\ 0.
\end{equation*}
Therefore, we can conclude that $\left|\frac{\log a_d(n)}{n^d}-\log\kappa_d\right|\ \le\ \frac{d\log2}n$, which is exactly item 2.
\end{proof}

\end{document}
